% slides.tex -- Beamer deck for the final review (MASTER_PLAN lead decision L5).
%
% It follows docs/PRESENTATION.md slide by slide; the speaker notes there are
% the long form of the \note{} text here. Rules, as for the paper:
%   - every number on a slide is a macro from paper/numbers.tex, which
%     paper/numbers.py writes from results/ (run `python reproduce.py paper`
%     first); none is typed by hand, and tools/check_doc_numbers.py lints
%     this file for hand-typed figures;
%   - every figure is a PDF from paper/figures/, by relative path.
% No LaTeX is installed locally (L6): compile on Overleaf with pdfLaTeX, the
% main document set to presentation/slides.tex (presentation/README.md).
%
% \pnum{name} prints the numbers.tex macro \name, or a bold ?? if it is not
% defined. The headline macros (\diskAdaptiveVsSeqCI, \ldVsVcfCI, ...) are
% always defined by numbers.tex and are used directly; \pnum guards the
% systematic families, whose names depend on which rows exist.

\documentclass[aspectratio=169,11pt]{beamer}
\usetheme{default}
\usecolortheme{default}
\setbeamertemplate{navigation symbols}{}
\setbeamertemplate{footline}[frame number]
\setbeamerfont{note page}{size=\small}
% Uncomment to print the speaker notes after each slide:
% \setbeameroption{show notes}

\usepackage[T1]{fontenc}
\usepackage{graphicx}

% Overleaf compiles from the project root or from this file's folder,
% depending on the project layout; accept both.
\graphicspath{{../paper/figures/}{paper/figures/}}
\IfFileExists{../paper/numbers.tex}{% numbers.tex -- GENERATED by paper/numbers.py from results.
% Do not edit: rerun `python reproduce.py paper`. Every number the prose
% quotes is one of these macros (MASTER_PLAN section 5; paper/check.py).
% MISSING input: disk/corpus_profile.csv
% MISSING input: disk/deltas.csv
% MISSING input: disk/kappa.csv
% MISSING input: disk/table_by_corpus.csv
% MISSING input: disk/table_main.csv
% MISSING input: gpu/disk_learned_depth.csv
% MISSING input: live/live_metrics.csv
% MISSING input: live/live_parity.csv
% MISSING input: realio/realio_summary.csv
% MISSING input: web/comparisons.csv
% MISSING input: web/dataset_profile.json
% MISSING input: web/decomposition.csv
% MISSING input: web/freshness.csv
% MISSING input: web/policy_metrics.csv
% MISSING input: web/seq_metrics.csv
\providecommand{\missingnumber}{\textbf{??}}
\providecommand{\fixturemark}[1]{\textcolor{red}{#1\textsuperscript{FIXTURE}}}
\providecommand{\paperFixtureData}{false}
% --- disk/corpus_profile.csv
\newcommand{\diskCorpora}{\missingnumber}
\newcommand{\diskTraces}{\missingnumber}
% --- disk/deltas.csv
\newcommand{\diskAdaptiveVsSeq}{\missingnumber}
\newcommand{\diskAdaptiveVsSeqCI}{\missingnumber}
\newcommand{\diskAdaptiveVsSeqHi}{\missingnumber}
\newcommand{\diskAdaptiveVsSeqLo}{\missingnumber}
\newcommand{\diskAdaptiveVsSeqLosses}{\missingnumber}
\newcommand{\diskAdaptiveVsSeqN}{\missingnumber}
\newcommand{\diskAdaptiveVsSeqWins}{\missingnumber}
\newcommand{\diskDeepEightIORatio}{\missingnumber}
\newcommand{\diskDeepEightIORatioCI}{\missingnumber}
\newcommand{\diskDeepEightIORatioHi}{\missingnumber}
\newcommand{\diskDeepEightIORatioLo}{\missingnumber}
\newcommand{\diskDeepEightIORatioLosses}{\missingnumber}
\newcommand{\diskDeepEightIORatioN}{\missingnumber}
\newcommand{\diskDeepEightIORatioWins}{\missingnumber}
\newcommand{\diskDeepEightVsNone}{\missingnumber}
\newcommand{\diskDeepEightVsNoneCI}{\missingnumber}
\newcommand{\diskDeepEightVsNoneHi}{\missingnumber}
\newcommand{\diskDeepEightVsNoneLo}{\missingnumber}
\newcommand{\diskDeepEightVsNoneLosses}{\missingnumber}
\newcommand{\diskDeepEightVsNoneN}{\missingnumber}
\newcommand{\diskDeepEightVsNoneWins}{\missingnumber}
\newcommand{\diskGuardVsSeq}{\missingnumber}
\newcommand{\diskGuardVsSeqCI}{\missingnumber}
\newcommand{\diskGuardVsSeqHi}{\missingnumber}
\newcommand{\diskGuardVsSeqLo}{\missingnumber}
\newcommand{\diskGuardVsSeqLosses}{\missingnumber}
\newcommand{\diskGuardVsSeqN}{\missingnumber}
\newcommand{\diskGuardVsSeqWins}{\missingnumber}
\newcommand{\diskLinuxRaIORatio}{\missingnumber}
\newcommand{\diskLinuxRaIORatioCI}{\missingnumber}
\newcommand{\diskLinuxRaIORatioHi}{\missingnumber}
\newcommand{\diskLinuxRaIORatioLo}{\missingnumber}
\newcommand{\diskLinuxRaIORatioLosses}{\missingnumber}
\newcommand{\diskLinuxRaIORatioN}{\missingnumber}
\newcommand{\diskLinuxRaIORatioWins}{\missingnumber}
\newcommand{\diskLinuxRaVsNone}{\missingnumber}
\newcommand{\diskLinuxRaVsNoneCI}{\missingnumber}
\newcommand{\diskLinuxRaVsNoneHi}{\missingnumber}
\newcommand{\diskLinuxRaVsNoneLo}{\missingnumber}
\newcommand{\diskLinuxRaVsNoneLosses}{\missingnumber}
\newcommand{\diskLinuxRaVsNoneN}{\missingnumber}
\newcommand{\diskLinuxRaVsNoneWins}{\missingnumber}
% --- gpu/disk_learned_depth.csv
\newcommand{\ldVsVcf}{\missingnumber}
\newcommand{\ldVsVcfCI}{\missingnumber}
\newcommand{\ldVsVcfHi}{\missingnumber}
\newcommand{\ldVsVcfLo}{\missingnumber}
\newcommand{\ldVsVcfLosses}{\missingnumber}
\newcommand{\ldVsVcfN}{\missingnumber}
\newcommand{\ldVsVcfWins}{\missingnumber}
% --- live/live_metrics.csv
\newcommand{\liveInvalidPasses}{\missingnumber}
% --- live/live_parity.csv
\newcommand{\liveParityMismatches}{\missingnumber}
% --- prefetch_common.utility
\newcommand{\diskAssumedC}{$0.526$}
\newcommand{\diskBreakEvenC}{$1.0101$}
% --- web/comparisons.csv
\newcommand{\webControllerVsBestFixed}{\missingnumber}
\newcommand{\webControllerVsBestFixedCI}{\missingnumber}
\newcommand{\webControllerVsBestFixedHi}{\missingnumber}
\newcommand{\webControllerVsBestFixedLo}{\missingnumber}
% --- web/decomposition.csv
\newcommand{\webFreshnessLever}{\missingnumber}
\newcommand{\webFreshnessLeverCI}{\missingnumber}
\newcommand{\webFreshnessLeverHi}{\missingnumber}
\newcommand{\webFreshnessLeverLo}{\missingnumber}
}{%
  \IfFileExists{paper/numbers.tex}{% numbers.tex -- GENERATED by paper/numbers.py from results.
% Do not edit: rerun `python reproduce.py paper`. Every number the prose
% quotes is one of these macros (MASTER_PLAN section 5; paper/check.py).
% MISSING input: disk/corpus_profile.csv
% MISSING input: disk/deltas.csv
% MISSING input: disk/kappa.csv
% MISSING input: disk/table_by_corpus.csv
% MISSING input: disk/table_main.csv
% MISSING input: gpu/disk_learned_depth.csv
% MISSING input: live/live_metrics.csv
% MISSING input: live/live_parity.csv
% MISSING input: realio/realio_summary.csv
% MISSING input: web/comparisons.csv
% MISSING input: web/dataset_profile.json
% MISSING input: web/decomposition.csv
% MISSING input: web/freshness.csv
% MISSING input: web/policy_metrics.csv
% MISSING input: web/seq_metrics.csv
\providecommand{\missingnumber}{\textbf{??}}
\providecommand{\fixturemark}[1]{\textcolor{red}{#1\textsuperscript{FIXTURE}}}
\providecommand{\paperFixtureData}{false}
% --- disk/corpus_profile.csv
\newcommand{\diskCorpora}{\missingnumber}
\newcommand{\diskTraces}{\missingnumber}
% --- disk/deltas.csv
\newcommand{\diskAdaptiveVsSeq}{\missingnumber}
\newcommand{\diskAdaptiveVsSeqCI}{\missingnumber}
\newcommand{\diskAdaptiveVsSeqHi}{\missingnumber}
\newcommand{\diskAdaptiveVsSeqLo}{\missingnumber}
\newcommand{\diskAdaptiveVsSeqLosses}{\missingnumber}
\newcommand{\diskAdaptiveVsSeqN}{\missingnumber}
\newcommand{\diskAdaptiveVsSeqWins}{\missingnumber}
\newcommand{\diskDeepEightIORatio}{\missingnumber}
\newcommand{\diskDeepEightIORatioCI}{\missingnumber}
\newcommand{\diskDeepEightIORatioHi}{\missingnumber}
\newcommand{\diskDeepEightIORatioLo}{\missingnumber}
\newcommand{\diskDeepEightIORatioLosses}{\missingnumber}
\newcommand{\diskDeepEightIORatioN}{\missingnumber}
\newcommand{\diskDeepEightIORatioWins}{\missingnumber}
\newcommand{\diskDeepEightVsNone}{\missingnumber}
\newcommand{\diskDeepEightVsNoneCI}{\missingnumber}
\newcommand{\diskDeepEightVsNoneHi}{\missingnumber}
\newcommand{\diskDeepEightVsNoneLo}{\missingnumber}
\newcommand{\diskDeepEightVsNoneLosses}{\missingnumber}
\newcommand{\diskDeepEightVsNoneN}{\missingnumber}
\newcommand{\diskDeepEightVsNoneWins}{\missingnumber}
\newcommand{\diskGuardVsSeq}{\missingnumber}
\newcommand{\diskGuardVsSeqCI}{\missingnumber}
\newcommand{\diskGuardVsSeqHi}{\missingnumber}
\newcommand{\diskGuardVsSeqLo}{\missingnumber}
\newcommand{\diskGuardVsSeqLosses}{\missingnumber}
\newcommand{\diskGuardVsSeqN}{\missingnumber}
\newcommand{\diskGuardVsSeqWins}{\missingnumber}
\newcommand{\diskLinuxRaIORatio}{\missingnumber}
\newcommand{\diskLinuxRaIORatioCI}{\missingnumber}
\newcommand{\diskLinuxRaIORatioHi}{\missingnumber}
\newcommand{\diskLinuxRaIORatioLo}{\missingnumber}
\newcommand{\diskLinuxRaIORatioLosses}{\missingnumber}
\newcommand{\diskLinuxRaIORatioN}{\missingnumber}
\newcommand{\diskLinuxRaIORatioWins}{\missingnumber}
\newcommand{\diskLinuxRaVsNone}{\missingnumber}
\newcommand{\diskLinuxRaVsNoneCI}{\missingnumber}
\newcommand{\diskLinuxRaVsNoneHi}{\missingnumber}
\newcommand{\diskLinuxRaVsNoneLo}{\missingnumber}
\newcommand{\diskLinuxRaVsNoneLosses}{\missingnumber}
\newcommand{\diskLinuxRaVsNoneN}{\missingnumber}
\newcommand{\diskLinuxRaVsNoneWins}{\missingnumber}
% --- gpu/disk_learned_depth.csv
\newcommand{\ldVsVcf}{\missingnumber}
\newcommand{\ldVsVcfCI}{\missingnumber}
\newcommand{\ldVsVcfHi}{\missingnumber}
\newcommand{\ldVsVcfLo}{\missingnumber}
\newcommand{\ldVsVcfLosses}{\missingnumber}
\newcommand{\ldVsVcfN}{\missingnumber}
\newcommand{\ldVsVcfWins}{\missingnumber}
% --- live/live_metrics.csv
\newcommand{\liveInvalidPasses}{\missingnumber}
% --- live/live_parity.csv
\newcommand{\liveParityMismatches}{\missingnumber}
% --- prefetch_common.utility
\newcommand{\diskAssumedC}{$0.526$}
\newcommand{\diskBreakEvenC}{$1.0101$}
% --- web/comparisons.csv
\newcommand{\webControllerVsBestFixed}{\missingnumber}
\newcommand{\webControllerVsBestFixedCI}{\missingnumber}
\newcommand{\webControllerVsBestFixedHi}{\missingnumber}
\newcommand{\webControllerVsBestFixedLo}{\missingnumber}
% --- web/decomposition.csv
\newcommand{\webFreshnessLever}{\missingnumber}
\newcommand{\webFreshnessLeverCI}{\missingnumber}
\newcommand{\webFreshnessLeverHi}{\missingnumber}
\newcommand{\webFreshnessLeverLo}{\missingnumber}
}{%
    \typeout{paper/numbers.tex not found: run python reproduce.py paper}}}
\providecommand{\missingnumber}{\textbf{??}}
\newcommand{\pnum}[1]{\ifcsname #1\endcsname\csname #1\endcsname\else\textbf{??}\fi}
\newcommand{\fig}[2][0.62\textheight]{\centering\includegraphics[height=#1,width=\linewidth,keepaspectratio]{#2}}

\title{What Makes Prefetching Pay?}
\subtitle{Disk blocks and product pages, simulated and measured}
\author{Ashvik Mishra (24BAI0342) \and Shourya Saraf (24BDS0391)}
\institute{School of Computer Science and Engineering\\
Vellore Institute of Technology, Vellore, India}
\date{Operating Systems course project}

\begin{document}

% ---------------------------------------------------------------- 1
\begin{frame}
  \titlepage
  \note{Prefetching fetches data before it is asked for: kernel read-ahead
  for files, browsers and proxies for web pages. Learning is often proposed
  to predict what to fetch. We tested, on public data and on real hardware,
  whether the prediction is what matters.}
\end{frame}

% ---------------------------------------------------------------- 2
\begin{frame}{The question, fixed before the experiments}
  \begin{columns}[T]
    \begin{column}{0.5\textwidth}
      Four levers of a prefetcher:
      \begin{enumerate}
        \item \textbf{how much} to prefetch (depth, or a growing window)
        \item \textbf{feedback} on its observed precision
        \item \textbf{freshness} of the statistics (web)
        \item the \textbf{learned predictor}
      \end{enumerate}
      \medskip
      Hypothesis: the predictor is the smallest lever.\\
      Comparisons and decision rules were pre-registered in
      \texttt{results/*/PREREGISTERED.md}.
    \end{column}
    \begin{column}{0.5\textwidth}
      \fig[0.6\textheight]{f1_schematic.pdf}
    \end{column}
  \end{columns}
  \note{Every claim follows the pre-registered rules: better, worse, or no
  difference detected, and when one lever may be called smaller than another.
  Including the claims that went against our hypothesis.}
\end{frame}

% ---------------------------------------------------------------- 3
\begin{frame}{One model for both studies}
  \begin{columns}[T]
    \begin{column}{0.48\textwidth}
      \begin{itemize}
        \item a cache of fixed capacity, a demand stream, a policy that
              issues speculative fetches
        \item hit ratio, precision, issued prefetches
        \item one objective:
              \[ U(c) = \frac{\mathit{hits} - c\cdot\mathit{issued}}{\mathit{requests}} \]
              with $c$ the cost of a prefetch in hits
        \item the classic disk cost model is $U$ at a $c$ above one: an
              artefact of charging every block as a device read
      \end{itemize}
    \end{column}
    \begin{column}{0.52\textwidth}
      \fig[0.55\textheight]{f8_utility_envelopes.pdf}
    \end{column}
  \end{columns}
  \note{The disk cost model and the web cost per prefetch are both U at some
  c, so one formula ranks policies in both domains. Definitions are in
  docs/METRICS.md.}
\end{frame}

% ---------------------------------------------------------------- 4
\begin{frame}{Disk data: real block traces, reconstructed carefully}
  \begin{columns}[T]
    \begin{column}{0.45\textwidth}
      \begin{itemize}
        \item \diskTraces{} traces from \diskCorpora{} corpora: MSR Cambridge,
              FIU, Alibaba, Tencent (libCacheSim \texttt{oracleGeneral})
        \item offsets in bytes; FIU records are pages, merged back into
              requests
        \item no read/write flag, no service times: writes replay as reads,
              and every replay names its cost model
      \end{itemize}
    \end{column}
    \begin{column}{0.55\textwidth}
      \fig[0.6\textheight]{f2_hit_vs_depth.pdf}
    \end{column}
  \end{columns}
  \note{Our own early notes described the format as one record per block. It
  is one record per request (per page for FIU). The loader detects units per
  file and every result names its reconstruction rule (DECISIONS D20).}
\end{frame}

% ---------------------------------------------------------------- 5
\begin{frame}{The disk simulator and its policies}
  \begin{columns}[T]
    \begin{column}{0.55\textwidth}
      Causal loop: demand $\rightarrow$ propose $\rightarrow$ prefetch;
      the first window never prefetches.
      \begin{itemize}
        \item fixed: \texttt{none}, \texttt{sequential}, \texttt{deep@N}
        \item per request: \texttt{stride}, \texttt{markov}, \texttt{lstm}
        \item trained classifier router: \texttt{adaptive}
        \item feedback gate: \texttt{guard}, \texttt{depth\_adaptive},
              \texttt{correlate}
        \item Linux read-ahead policy: \texttt{linux\_ra@N}
        \item learned depth: \texttt{learned\_depth}
      \end{itemize}
    \end{column}
    \begin{column}{0.45\textwidth}
      \fig[0.55\textheight]{f1_schematic.pdf}
    \end{column}
  \end{columns}
  \note{Same requests, cache and warm-up for every policy. For the Linux
  read-ahead policy and the learned depth model the tests replace the future
  with random requests and get every prefix outcome unchanged; the streamed
  replay equals the in-memory one bit for bit.}
\end{frame}

% ---------------------------------------------------------------- 6
\begin{frame}{Disk result 1: how much to prefetch dominates}
  \begin{columns}[T]
    \begin{column}{0.42\textwidth}
      Mean block hit ratio, 1\,MiB cache:
      \begin{itemize}
        \item \texttt{none}: \pnum{diskHitNoneOneMiB}
        \item \texttt{deep@8} vs \texttt{none}: \diskDeepEightVsNoneCI
        \item \texttt{deep@256} vs \texttt{none}:
              \pnum{diskDeltaHitRatioDeepAtTwoHundredFiftySixVsNoneOneMiBCI}
        \item \texttt{linux\_ra@256} vs \texttt{none}: \diskLinuxRaVsNoneCI
        \item Belady MIN (demand only, \emph{not} a prefetch bound):
              \pnum{diskHitBeladyMinOneMiB}
      \end{itemize}
    \end{column}
    \begin{column}{0.58\textwidth}
      \fig{f2_hit_vs_depth.pdf}
    \end{column}
  \end{columns}
  \note{Hit ratio keeps climbing with depth well past 8, and does not
  saturate with cache size (D20 corrects both earlier claims). Belady MIN is
  optimal demand replacement without prefetching; deep read-ahead exceeds it.}
\end{frame}

% ---------------------------------------------------------------- 7
\begin{frame}{Disk result 2: hit ratio is not device work}
  \begin{columns}[T]
    \begin{column}{0.42\textwidth}
      Device I/Os relative to \texttt{none}, 1\,MiB:
      \begin{itemize}
        \item \texttt{deep@8}: \diskDeepEightIORatioCI
        \item \texttt{linux\_ra@256}: \diskLinuxRaIORatioCI
      \end{itemize}
      Blocks moved:
      \begin{itemize}
        \item \texttt{deep@256}:
              \pnum{diskDeltaBlocksMovedDeepAtTwoHundredFiftySixVsNoneOneMiBCI}
        \item \texttt{linux\_ra@256}:
              \pnum{diskDeltaBlocksMovedLinuxRaAtTwoHundredFiftySixVsNoneOneMiBCI}
      \end{itemize}
    \end{column}
    \begin{column}{0.58\textwidth}
      \fig{f3_device_io_vs_hit.pdf}
    \end{column}
  \end{columns}
  \note{Fixed-depth read-ahead raises hit ratio while issuing more device I/Os
  than doing nothing; only the growing Linux-style window raises hit ratio and
  lowers I/Os. Every hit ratio is shown with these two ratios.}
\end{frame}

% ---------------------------------------------------------------- 8
\begin{frame}{Disk result 3: the learned models}
  \begin{columns}[T]
    \begin{column}{0.5\textwidth}
      Hit ratio against \texttt{sequential}:
      \begin{itemize}
        \item classifier router, 1\,MiB: \diskAdaptiveVsSeqCI
        \item classifier router, 64\,MiB:
              \pnum{diskDeltaHitRatioAdaptiveVsSequentialSixtyFourMiBCI}
        \item feedback gate, 1\,MiB: \diskGuardVsSeqCI
      \end{itemize}
      Learned depth against a validation-chosen fixed depth
      (\ldVsVcfN{} traces, \pnum{ldSeedsCapTwoHundredFiftySix} seeds): \ldVsVcfCI; it won on \ldVsVcfWins.
    \end{column}
    \begin{column}{0.5\textwidth}
      \fig[0.55\textheight]{f3_device_io_vs_hit.pdf}
    \end{column}
  \end{columns}
  \note{The trained router loses to two-block read-ahead. The GPU depth model
  loses to the strongest fair baseline. The predictor is smaller than depth
  and window growth but not the smallest lever: at 64 MiB the feedback gate's
  effect is smaller still (D20).}
\end{frame}

% ---------------------------------------------------------------- 9
\begin{frame}{When does a policy pay on a device? $\kappa^*$}
  \begin{columns}[T]
    \begin{column}{0.42\textwidth}
      \begin{itemize}
        \item device time $T = N_{io}\,t_{pos} + N_{blk}\,t_{xfer}$
        \item measured $\kappa = t_{pos}/t_{xfer}$:
              HDD \pnum{deviceSdbKappa}, SSD \pnum{deviceSsdKappa} blocks
        \item $\kappa^*$: the device ratio above which a policy costs less
              device time than \texttt{none}
        \item more I/Os than \texttt{none}: it never pays
      \end{itemize}
    \end{column}
    \begin{column}{0.58\textwidth}
      \fig{f4_kappa_cdf.pdf}
    \end{column}
  \end{columns}
  \note{The vertical lines are the two drives we calibrated. At 1\,MiB the
  Linux-style window is cheaper than doing nothing on both, on almost every
  trace; fixed depth 8 on very few (DECISIONS D18).}
\end{frame}

% ---------------------------------------------------------------- 10
\begin{frame}{The real Linux testbed}
  \begin{columns}[T]
    \begin{column}{0.55\textwidth}
      The simulator decides; the testbed does the I/O.
      \begin{itemize}
        \item Backend B: user-space LRU, \texttt{O\_DIRECT}, asynchronous
              prefetch; hit/miss asserted equal to the simulator on every
              request
        \item Backend A: kernel page cache in a memory cgroup, chunked
              \texttt{POSIX\_FADV\_WILLNEED}; \texttt{kernel\_ra} measures
              Linux's own read-ahead
        \item fresh process per run, file dropped from the page cache,
              drive cache displaced, shuffled repeats, quiet host only
      \end{itemize}
    \end{column}
    \begin{column}{0.45\textwidth}
      \fig[0.55\textheight]{f5_sim_vs_measured.pdf}
    \end{column}
  \end{columns}
  \note{Kernel facts measured first: a WILLNEED hint is truncated to the
  device's request limit, so hints are chunked; pages charged to another
  cgroup are not re-charged, so the file is dropped before every run;
  cachestat plus a non-blocking read separate ready hits from late ones.}
\end{frame}

% ---------------------------------------------------------------- 11
\begin{frame}{Real device: the simulator's favourite was never the fastest}
  \begin{columns}[T]
    \begin{column}{0.42\textwidth}
      \begin{itemize}
        \item in no measured cell did the highest simulated hit ratio give
              the lowest measured latency
        \item doing nothing was often the fastest
        \item simulated hit ratio did not rank policies by measured latency
              on the HDD
      \end{itemize}
      Per trace, think model and device; never a property of all disks.
    \end{column}
    \begin{column}{0.58\textwidth}
      \fig{f5_sim_vs_measured.pdf}
    \end{column}
  \end{columns}
  \note{A prefetch still in flight when it is needed, or one that queues in
  front of a demand read, costs time the simulator never charges. Example and
  per-cell values: docs/PRESENTATION.md slide 11 and paper Table 3.}
\end{frame}

% ---------------------------------------------------------------- 12
\begin{frame}{Web study: RetailRocket, and the stale index}
  \begin{columns}[T]
    \begin{column}{0.42\textwidth}
      \begin{itemize}
        \item per-session browser cache, four chronological splits: index,
              ranker, validation, test
        \item every policy selected on validation \emph{at each cache size}
        \item freshness lever, causal online counts vs the frozen logistic
              ranker: \webFreshnessLeverCI
      \end{itemize}
    \end{column}
    \begin{column}{0.58\textwidth}
      \fig{f6_freshness.pdf}
    \end{column}
  \end{columns}
  \note{The original branch's candidate index came from the oldest data only;
  most of its "coverage ceiling" was staleness. Selecting at one cache size
  and reusing the choice hid losses to no prefetching (W1, W2).}
\end{frame}

% ---------------------------------------------------------------- 13
\begin{frame}{Web result: what moves the needle}
  \begin{columns}[T]
    \begin{column}{0.5\textwidth}
      Hit ratio, cache of 8 objects:
      \begin{itemize}
        \item depth:
              \pnum{webLeverDepthFixedFourFixedOneLogisticLogisticEightZeroPointZeroCI}
        \item ranker, counting vs logistic:
              \pnum{webLeverRankerFixedOneFixedOneTransSZeroLogisticEightZeroPointZeroCI}
        \item GRU vs online counting:
              \pnum{webLeverSequenceModelFixedFourFixedFourSeqTransOnlineEightZeroPointZeroCI}
      \end{itemize}
      Validated adaptive vs best fixed depth (utility):
      \webControllerVsBestFixedCI
    \end{column}
    \begin{column}{0.5\textwidth}
      \fig[0.55\textheight]{f6_freshness.pdf}
    \end{column}
  \end{columns}
  \note{On the web the learned ranker moves hit ratio less than depth or
  freshness, but the GPU sequence model's effect overlaps depth, so by the
  pre-registered rule the predictor is not called the smallest lever. The
  sequence model is worse than counting transitions. The adaptive controller still beats the best fixed depth at
  caches 8 and 32 for an assumed cost per prefetch.}
\end{frame}

% ---------------------------------------------------------------- 14
\begin{frame}{Live HTTP replay: measured waits, exact parity}
  \begin{columns}[T]
    \begin{column}{0.42\textwidth}
      \begin{itemize}
        \item sessions whose counters differ from the simulator:
              \liveParityMismatches; invalid passes: \liveInvalidPasses
        \item HDD store, online-index controller: saved
              \pnum{liveSavedHddAdaptiveZeroPointOneAtTransOnlineCI} per
              view at \pnum{liveBytesXHddAdaptiveZeroPointOneAtTransOnline}
              the bytes
        \item fixed depth 4: \pnum{liveSavedHddFixedFourCI} at
              \pnum{liveBytesXHddFixedFour}
      \end{itemize}
      Network emulated, object sizes assumed, waits measured.
    \end{column}
    \begin{column}{0.58\textwidth}
      \fig{f7_live_saved_vs_bytes.pdf}
    \end{column}
  \end{columns}
  \note{Each emulated browser's cache is the simulator's cache class, so the
  counts must match exactly, and they did. With a FIFO link (SSD store, one
  repetition) depth 8's saving was no longer detectable: speculative bytes
  queued in front of demand fetches (W4).}
\end{frame}

% ---------------------------------------------------------------- 15
\begin{frame}{Would a better learner help?}
  \begin{columns}[T]
    \begin{column}{0.46\textwidth}
      \small Pre-registered after E1--E6 were known.
      \begin{itemize}
        \item \pnum{clfSynTunedFamilies} tuned classifier families in the
              router: \pnum{decClfClaimBeatsReadAheadOneMiBVerdict}
              (vs \texttt{sequential}, 64\,MiB, best:
              \pnum{decClfOneMlpSixtyFourMiBPairedCI}); device I/Os
              \pnum{clfRouterIORatioMinOneMiB}--\pnum{clfRouterIORatioMaxOneMiB}
              of \texttt{none}'s
        \item depth as an action, hindsight labels: vs a fixed depth chosen per
              trace, \pnum{decHsrHOneHitRatioRoutedHitVsVcfHitOneMiBPairedCI}
              (\pnum{decHsrHOneHitRatioRoutedHitVsVcfHitOneMiBVerdict})
        \item web: tuned MLP ranker raises the controller's $U(c)$ by
              \pnum{decWrkSOneMlpUtilityPerViewCapEightCI}, still below online
              counts (\pnum{decWrkKFiveMlpHitRatioCapEightCI} hit ratio)
      \end{itemize}
    \end{column}
    \begin{column}{0.54\textwidth}
      \fig[0.6\textheight]{f9_action_space.pdf}
    \end{column}
  \end{columns}
  \note{The router's four classes each prefetch at most two blocks, so any
  classifier is capped near sequential; the classifiers disagree with GNB on
  most windows yet the hit ratio barely moves. Give a router depth as an
  action and it only matches a fixed depth chosen per trace. On the web a
  better ranker helps a little; freshness still decides (D22, W5).}
\end{frame}

% ---------------------------------------------------------------- 16
\begin{frame}{Across both domains}
  \begin{columns}[T]
    \begin{column}{0.42\textwidth}
      \begin{enumerate}
        \item how much to prefetch, and in what shape, mattered more than
              any learned predictor we tried
        \item a simulated hit ratio is neither latency nor device work
        \item on the web, freshness beat sophistication
      \end{enumerate}
      \medskip
      \small Limits: writes replayed as reads; one HDD and a shared SSD;
      emulated network; assumed web costs.
    \end{column}
    \begin{column}{0.58\textwidth}
      \fig{f8_utility_envelopes.pdf}
    \end{column}
  \end{columns}
  \note{Read the envelope left to right: at low cost per prefetch deep
  policies win; as the cost rises the best policy issues less, until doing
  nothing wins. Where a domain sits on that axis decides which policy pays.
  Every number is generated from results/, every comparison pre-registered,
  and the negative results are reported as found.}
\end{frame}

\end{document}
